\section*{Capítulo \ref{ch:b3:representations} --- Representaciones
de los grupos finitos}

\begin{solution}{exo:b3:representations:1}
(a) $\Z/n\Z$ es abeliano: todas las
\omterm{def:b3:representations:rep}{irreducibles} tienen grado $1$
(la~\cref{prop:b3:representations:onedim}), es decir, son morfismos
$\chi \colon \Z/n\Z \to \C^\times$, determinados por $\chi(\bar 1)
= \omega$ con $\omega^n = 1$: los $n$
\omterm{def:b3:representations:character}{caracteres} $\chi_k(\bar
m) = \eu^{2\iu\pi km/n}$. Para $n = 4$ (clases de $= $ elementos
$\bar0,\bar1,\bar2,\bar3$):
\[
\begin{array}{c|cccc}
& \bar0 & \bar1 & \bar2 & \bar3\\
\hline
\chi_0 & 1 & 1 & 1 & 1\\
\chi_1 & 1 & \iu & -1 & -\iu\\
\chi_2 & 1 & -1 & 1 & -1\\
\chi_3 & 1 & -\iu & -1 & \iu
\end{array}
\]
(b) Filas: $\langle\chi_k,\chi_l\rangle = \frac14\sum_m
\eu^{2\iu\pi(l-k)m/4} = \delta_{kl}$ (suma geométrica); las columnas,
análogamente. La matriz $(\eu^{2\iu\pi km/n})_{k,m}$ es la matriz de la
\emph{transformada de Fourier discreta} --- el mismo filtro de raíces de
la unidad que se usó en el capítulo de funciones generatrices de segundo
año; la \omterm{thm:b3:representations:orthogonality}{ortogonalidad de
los caracteres} generaliza la fórmula de inversión de la TFD.
\end{solution}

\begin{solution}{exo:b3:representations:2}
(a) La matriz de $\rho(g)$ en la base $(e_x)$ es una matriz de
permutación, cuya traza es el número de $x$ con $g\cdot x =
x$. Entonces
\[
\langle\mathbf 1, \chi\rangle = \frac1{\abs G}\sum_g
\abs{\operatorname{Fix}(g)} = \#\{\text{órbitas}\}
\]
por el lema de recuento de Burnside (el~\cref{exo:b3:groups:5}) ---
equivalentemente, esto calcula la multiplicidad de la
\omterm{def:b3:representations:rep}{representación} trivial, cuyo
espacio isotípico es el espacio de los vectores $G$-invariantes, de
dimensión igual al número de \omterm{def:b3:groups:action}{órbitas} (un
indicador por \omterm{def:b3:groups:action}{órbita}).

(b) Como $\operatorname{Fix}_{X\times X}(g) =
\operatorname{Fix}_X(g)^2$, el apartado (a) aplicado a $X \times X$ da
$\langle\chi,\chi\rangle = \frac1{\abs
G}\sum\abs{\operatorname{Fix}(g)}^2 = \#\{\text{órbitas en }
X\times X\}$ ($\chi$ es real). Escribiendo $\chi = \mathbf 1 +
\psi$: $\langle\mathbf1,\chi\rangle = 1$ (transitividad), luego
$\langle\psi,\psi\rangle = \langle\chi,\chi\rangle - 1$. La acción sobre
$X\times X$ tiene la diagonal como una
\omterm{def:b3:groups:action}{órbita}; hay exactamente otra
\omterm{def:b3:groups:action}{órbita} si y solo si $G$ es transitiva
sobre los pares de puntos distintos: $\langle\psi,\psi\rangle = 1$ si y
solo si es $2$-transitiva
(\cref{cor:b3:representations:multiplicity}).

(c) $S_n$ es $2$-transitivo sobre $\{1,\dots,n\}$ (todo par de puntos
distintos puede ir a cualquier otro): $\psi = \chi_{\mathrm{std}}$ es
\omterm{def:b3:representations:rep}{irreducible}.
\end{solution}

\begin{solution}{exo:b3:representations:3}
$S_3$: tres clases, $\sum n_i^2 = 6 = 1 + 1 + 4$; los dos
\omterm{def:b3:representations:character}{caracteres} de grado $1$ son
$\mathbf 1, \varepsilon$ ($[S_3 :
A_3] = 2$); la tercera fila $(2, a, b)$ se sigue de la ortogonalidad de
columnas con la columna de $e$: $1 - 1 + 2a = 0$ y $1 + 1 +
2b = 0$: $a = 0$, $b = -1$ --- la tabla de
\cref{ex:b3:representations:s3}.

$S_4$ comprobaciones (filas): $\langle\chi_{\mathrm{std}},
\varepsilon\chi_{\mathrm{std}}\rangle = \frac1{24}(9 - 6 + 3 + 0
- 6) = 0$; $\langle\chi_2,\chi_2\rangle = \frac1{24}(4 + 0 + 12
+ 8 + 0) = 1$. Columnas: $e$ contra $(1\,2)$: $1 - 1 + 0 + 3 - 3
= 0$; $(1\,2)$ consigo misma: $1 + 1 + 0 + 1 + 1 = 4 =
\abs{Z_{S_4}((1\,2))} = 24/6$.

\omterm{def:b3:representations:character}{Carácter} de permutación sobre
$4$ puntos: $(4, 2, 0, 1, 0) = \mathbf
1 + \chi_{\mathrm{std}}$ (recuentos de puntos fijos; réstese la fila
superior). Para $\chi_{\mathrm{std}}^2 = (9, 1, 1, 0, 1)$:
\[
\langle\mathbf1,\cdot\rangle = \tfrac{9 + 6 + 3 + 0 + 6}{24} =
1,\quad
\langle\varepsilon,\cdot\rangle = 0,\quad
\langle\chi_2,\cdot\rangle = \tfrac{18 + 6}{24} = 1,\quad
\langle\chi_{\mathrm{std}},\cdot\rangle = 1,\quad
\langle\varepsilon\chi_{\mathrm{std}},\cdot\rangle = 1:
\]
$\chi_{\mathrm{std}}^2 = \mathbf 1 + \chi_2 +
\chi_{\mathrm{std}} + \varepsilon\chi_{\mathrm{std}}$
(dimensiones: $9 = 1 + 2 + 3 + 3$).
\end{solution}

\begin{solution}{exo:b3:representations:4}
Ambos grupos tienen cinco clases y patrón de grados $(1,1,1,1,2)$
(cuatro de grado $1$ procedentes de la abelianización
$\cong (\Z/2\Z)^2$, y después $\sum n_i^2 = 8$). Ordenando las clases
$e$, $z$ (la involución central: $r^2$ y $-1$, respectivamente) y las
tres clases de dos elementos:
\[
\begin{array}{c|ccccc}
& e & z & C_1 & C_2 & C_3\\
\hline
\chi^{(1)} & 1 & 1 & 1 & 1 & 1\\
\chi^{(2)} & 1 & 1 & 1 & -1 & -1\\
\chi^{(3)} & 1 & 1 & -1 & 1 & -1\\
\chi^{(4)} & 1 & 1 & -1 & -1 & 1\\
\chi^{(5)} & 2 & -2 & 0 & 0 & 0
\end{array}
\]
(la última fila, por ortogonalidad de columnas). Tablas idénticas para
$D_4$ y $Q_8$, que no son isomorfos (el~\cref{pb:b3:groups:1}). La tabla
\emph{sí} capta: $\abs
G$, los cardinales de las clases, el
\omterm{ex:b3:groups:actions}{centro} ($\{g : \abs{\chi_i(g)} = n_i\
\forall i\}$: orden $2$ en ambos), la abelianización y todo el retículo
de \omterm{def:b3:groups:normal}{subgrupos normales} (núcleos e
intersecciones, el~\cref{exo:b3:representations:6}). No capta los
órdenes de los elementos: $D_4$ tiene cinco involuciones y $Q_8$ solo
una --- de modo que el tipo de isomorfía es genuinamente más fino que la
\omterm{def:b3:representations:table}{tabla de caracteres}.
\end{solution}

\begin{solution}{exo:b3:representations:5}
(a) En $S_4$, el centralizador de $(1\,2\,3)$ tiene orden $24/8 =
3$: es $\langle(1\,2\,3)\rangle \subseteq A_4$. Luego
$Z_{A_4}((1\,2\,3))$ tiene orden $3$ y la $A_4$-clase tiene $12/3
= 4$ elementos: los ocho $3$-ciclos se desdoblan en dos $A_4$-clases
(representadas por $(1\,2\,3)$ y su inverso).

(b) $A_4/V \cong \Z/3\Z$ da tres
\omterm{def:b3:representations:character}{caracteres} de grado $1$
($\omega = \eu^{2\iu\pi/3}$; las clases de $3$-ciclos van a
$\bar1, \bar2$); la restricción de $\chi_{\mathrm{std}}$ sigue siendo
\omterm{def:b3:representations:rep}{irreducible}
($\langle\chi,\chi\rangle = \frac1{12}(9 + 3
\cdot 1 + 0 + 0) = 1$):
\[
\begin{array}{c|cccc}
A_4 & e\,[1] & (1\,2)(3\,4)\,[3] & (1\,2\,3)\,[4] &
(1\,3\,2)\,[4]\\
\hline
\mathbf 1 & 1 & 1 & 1 & 1\\
\chi_\omega & 1 & 1 & \omega & \omega^2\\
\chi_{\bar\omega} & 1 & 1 & \omega^2 & \omega\\
\chi_3 & 3 & -1 & 0 & 0
\end{array}
\]
(c) Núcleos: $\ker\chi_\omega = \ker\chi_{\bar\omega} = V$;
$\ker\chi_3 = \{e\}$ (ninguna otra entrada tiene módulo $3$). Los
\omterm{def:b3:groups:normal}{subgrupos normales} son las intersecciones
de núcleos (el~\cref{exo:b3:representations:6}): $\{e\}$, $V$, $A_4$ ---
en particular, $A_4$ no tiene \omterm{def:b3:groups:normal}{subgrupo
normal} de orden $2$ ni de índice
$2$.
\end{solution}

\begin{solution}{exo:b3:representations:6}
(a) Si $\chi(g) = \chi(e) = n$: la igualdad en $\abs{\chi(g)} \leq n$
obliga a $\rho(g) = \lambda\,\mathrm{id}$
(la~\cref{prop:b3:representations:charbasics}) con $n\lambda = n$:
$\rho(g) = \mathrm{id}$. El recíproco es claro.

(b) Sea $N \trianglelefteq G$. La
\omterm{ex:b3:representations:examples}{representación regular} de $G/N$
es fiel; descompóngase en
\omterm{def:b3:representations:rep}{irreducibles} de $G/N$ y elévense a
$G$ (la~\cref{prop:b3:representations:onedim}(b)): se obtienen
\omterm{def:b3:representations:character}{caracteres}
\omterm{def:b3:representations:rep}{irreducibles} $\chi_{i_1}, \dots$ de
$G$ cuyos núcleos contienen $N$ y cuyo núcleo \emph{común} es
exactamente la preimagen de $\{e\}$, es decir, $N$ (fidelidad sobre el
cociente). Luego $N = \bigcap_j \ker\chi_{i_j}$.

(c) Si $G$ es simple: para un
\omterm{def:b3:representations:rep}{irreducible} no trivial $\chi$,
$\ker\chi \trianglelefteq G$ no es $G$ (una
\omterm{def:b3:representations:rep}{representación irreducible} trivial
sobre todo $G$ es el \omterm{def:b3:representations:character}{carácter}
trivial), luego $\ker\chi = \{e\}$. Recíprocamente, supóngase que todos
los núcleos no triviales son triviales y sea $N \trianglelefteq G$ con
$N \neq
G$. En la expresión de $N$ como intersección de núcleos de (b), alguno
de los \omterm{def:b3:representations:character}{caracteres} implicados
es no trivial (si todos fueran triviales, la intersección sería $G$), y
su núcleo es $\{e\}$: $N =
\{e\}$. Así pues, los únicos \omterm{def:b3:groups:normal}{subgrupos
normales} son $\{e\}$ y $G$.
\end{solution}

\begin{solution}{exo:b3:representations:7}
Orden $8$ no abeliano: el número de
\omterm{def:b3:representations:character}{caracteres} de grado $1$ es
$[G : D(G)]$, un divisor propio de $8$ (no abeliano: $D(G) \neq
\{e\}$), y $\sum n_i^2 = 8$. Con $k$ unos y los grados restantes
$\geq 2$: $8 - k \equiv 0$ con cuadrados $\geq 4$, y $k
\mid 8$, $k < 8$. $k = 4$: un grado $2$ --- coherente. $k =
2$: quedan $6$, que no es suma de cuadrados $\geq 4$. $k = 1$:
imposible, pues $k = [G : D(G)] \geq 2$ --- $G/D(G)$ es un $2$-grupo
abeliano no trivial, ya que $G$ es un $2$-grupo con $D(G) \neq G$ por la
\omterm{def:b3:groups:derived}{resolubilidad} de los $p$-grupos
(el~\cref{ex:b3:groups:solvableexamples}). Así pues, el patrón es
$(1,1,1,1,2)$. Fidelidad de $\chi_5$: los cuatro
\omterm{def:b3:representations:character}{caracteres} de grado $1$
contienen todos $D(G)$ en sus núcleos; si $\ker\chi_5
\supseteq N \neq \{e\}$ para algún normal minimal $N$ ($N
\subseteq D(G)$ o no --- tómese $N \subseteq \ker\chi_5$ no trivial),
entonces $N$ estaría en los cinco núcleos, cuya intersección es trivial
(la \omterm{ex:b3:representations:examples}{representación regular} es
fiel): contradicción. Orden $p^3$ no abeliano: $Z(G)$ tiene orden $p$
(el orden $p^2$ haría $G/Z$ cíclico y $G$ abeliano), $G/Z(G)$, de orden
$p^2$, es abeliano, luego $D(G) \subseteq Z(G)$ y $D(G) \ne \{e\}$:
$D(G) = Z(G)$, lo que da $p^2$
\omterm{def:b3:representations:character}{caracteres} de grado $1$. Los
grados restantes cumplen $\sum n_i^2 = p^3 -
p^2$ con cada $n_i > 1$ divisor de $\abs G$
(el~\cref{pb:b3:representations:1}, pregunta 8), luego $n_i \in
\{p\}$ ($n_i = p^2$ se pasaría: $p^4 > p^3 - p^2$): exactamente $p
- 1$ \omterm{def:b3:representations:character}{caracteres} de grado $p$.
\end{solution}

\begin{solution}{exo:b3:representations:8}
Defínase, para representaciones
$\rho, \sigma$ de $G, H$ sobre $V, W$, la
\omterm{def:b3:representations:rep}{representación}
$\rho \boxtimes \sigma$ de $G \times H$ sobre $V \otimes W$ ---
concretamente, sobre matrices: $(\rho\boxtimes
\sigma)(g,h)$ es el producto de Kronecker $\rho(g)\otimes
\sigma(h)$, cuya traza es
$\operatorname{tr}\rho(g)\operatorname{tr}\sigma(h) =
\chi(g)\psi(h)$ (el producto de Kronecker de las matrices $A \otimes B$
tiene traza $\operatorname{tr}A\operatorname{tr}B$: su diagonal es
$a_{kk}b_{ll}$). Luego $\chi\psi$ es un
\omterm{def:b3:representations:character}{carácter}, y
\[
\langle\chi\psi, \chi'\psi'\rangle_{G\times H}
= \frac{1}{\abs G\abs H}\sum_{g,h}
\overline{\chi(g)\psi(h)}\,\chi'(g)\psi'(h)
= \langle\chi,\chi'\rangle_G\,\langle\psi,\psi'\rangle_H
= \delta_{\chi\chi'}\delta_{\psi\psi'}.
\]
En particular $\langle\chi\psi,\chi\psi\rangle = 1$: cada $\chi\psi$ es
\omterm{def:b3:representations:rep}{irreducible}
(el~\cref{cor:b3:representations:multiplicity}). Estos son $r_Gr_H$
\omterm{def:b3:representations:character}{caracteres}
\omterm{def:b3:representations:rep}{irreducibles} distintos; las clases
de $G\times H$ son los productos de clases ($(g,h) \sim (g',h')$
componente a componente), de modo que hay $r_Gr_H$: la lista es completa
(el~\cref{thm:b3:representations:numberirr}). Para $(\Z/2\Z)^2$: los
cuatro \omterm{def:b3:representations:character}{caracteres} de signos
$(\pm1)\otimes(\pm1)$ --- el bloque superior izquierdo de la tabla
del~\cref{exo:b3:representations:4}. Un grupo abeliano finito es
producto de grupos cíclicos (el~\cref{cor:b3:modules:abelian}); sus
\omterm{def:b3:representations:character}{caracteres}
\omterm{def:b3:representations:rep}{irreducibles} son productos de los
cíclicos (el~\cref{exo:b3:representations:1}): todos de grado $1$.
\end{solution}

\begin{solution}{exo:b3:representations:9}
(a) $D(A_5) = A_5$ (simple no abeliano), de modo que el único
\omterm{def:b3:representations:character}{carácter} de grado $1$ es
$\mathbf 1$ (la~\cref{prop:b3:representations:onedim}). Necesitamos
$n_2^2 + n_3^2
+ n_4^2 + n_5^2 = 59$ con cada $n_i \geq 2$; probando los cuadrados
$4, 9, 16, 25, 36, 49$: el único multiconjunto que funciona es
$\{9, 9, 16, 25\}$; con mayor cuadrado $49$, el resto $10$ no es suma de
tres cuadrados $\geq 4$; con $36$, el resto $23$ tampoco lo es
($16 + 4 + 4 = 24$, $9 + 9 + 4 = 22$); con mayor $25$ se comprueba que
$25 + 16 + 9 + 9 = 59$ funciona y que las variantes $25 +
25$, $25 + 16 + 16$, $25 + 16 + 4$ fallan; con mayor $16$:
$16\cdot3 = 48 < 59 - 4$. Grados: $1, 3, 3, 4, 5$.

(b) Permutación sobre $5$ puntos: puntos fijos $(5, 1, 2, 0, 0)$, luego
$\chi_4 = (4, 0, 1, -1, -1)$ con
$\langle\chi_4,\chi_4\rangle = \frac{16 + 0 + 20 + 12 + 12}{60}
= 1$: \omterm{def:b3:representations:rep}{irreducible}. Acción sobre los
seis $5$-subgrupos de Sylow: una involución fija exactamente $2$ (los
\omterm{ex:b3:groups:actions}{normalizadores} son diédricos de orden
$10$ y cada uno contiene $5$ involuciones: $30$ incidencias para $15$
involuciones), un elemento de orden $3$ fija $0$ (no hay orden $3$ en
$D_5$) y un elemento de orden $5$ fija exactamente $1$ (está en un único
Sylow): \omterm{def:b3:representations:character}{carácter} de
permutación $(6, 2, 0, 1, 1)$, y $\chi_5 =
(5, 1, -1, 0, 0)$ con norma $\frac{25 + 15 + 20 + 0 + 0}{60} =
1$: \omterm{def:b3:representations:rep}{irreducible}. Quedan dos filas
$(3, a, b, c, d)$, $(3, a', b', c',
d')$. Normas de columnas (el~\cref{cor:b3:representations:column}): en
la clase de
$(1\,2)(3\,4)$, $\abs{Z} = 4$: $1 + 0 + 1 + a^2 + a'^2 = 4$;
columna contra $e$: $1\cdot1 + 4\cdot 0 + 5\cdot1 + 3(a + a') =
0$: $a + a' = -2$, $a^2 + a'^2 = 2$: $a = a' = -1$. Sobre los
$3$-ciclos, $\abs Z = 3$: $1 + 1 + 1 + b^2 + b'^2 = 3$: $b = b'
= 0$. En cada $5$-clase, $\abs Z = 5$: la columna contra $e$ da
$1\cdot 1 + 4\cdot(-1) + 5\cdot 0 + 3(c + c') = 0$, luego $c + c' = 1$;
y $c^2 + c'^2 = 3$: $\{c, c'\} = \bigl\{\frac{1+\sqrt5}2,
\frac{1-\sqrt5}2\bigr\}$ --- la razón áurea y su conjugada; la segunda
$5$-clase lleva los valores intercambiados (las dos filas han de ser
ortogonales).

(c) En la tabla terminada, ninguna entrada de una fila no trivial
coincide con su grado fuera de la primera columna: todo núcleo $\{g :
\chi_i(g) = n_i\}$ es trivial. Por
el~\cref{exo:b3:representations:6}(c), $A_5$ es simple.
\end{solution}

\begin{solution}{exo:b3:representations:10}
$\rho(z)$ conmuta con todo $\rho(g)$ ($z$ es central), es decir,
$\rho(z) \in \operatorname{End}_G(V) = \C\,\mathrm{id}$ (Schur):
$\rho(z) = \lambda_z\,\mathrm{id}$, y $z \mapsto \lambda_z$ es
multiplicativa: un morfismo $Z(G) \to \C^\times$. Entonces
$\chi(z) = n\lambda_z$ con $\abs{\lambda_z} = 1$ (raíz de la unidad):
$\abs{\chi(z)} = n$. Si $\rho$ es fiel, $\lambda$ es inyectivo sobre
$Z(G)$ ($\rho(z) = \mathrm{id} \iff \lambda_z =
1$), luego $Z(G)$ se sumerge en $\C^\times$; y un subgrupo finito del
grupo multiplicativo de un cuerpo es cíclico
(\cref{thm:b3:galois:cyclic}).
\end{solution}

\begin{solution}{exo:b3:representations:11}
(a) Equivarianza: $\rho(h)p_i\rho(h)^{-1}$ reindexa la suma
($g \mapsto hgh^{-1}$, y $\chi_i$ es una
\omterm{def:b3:representations:character}{función de clase}): $p_i$
conmuta con la acción. Sobre una
\omterm{def:b3:representations:rep}{irreducible} $W$ de
\omterm{def:b3:representations:character}{carácter} $\chi_j$, Schur hace
de la restricción un escalar $\lambda\,\mathrm{id}_W$; tomando trazas,
\[
\lambda\,n_j = \frac{n_i}{\abs G}\sum_g
\overline{\chi_i(g)}\,\chi_j(g) = n_i\,\langle\chi_i,
\chi_j\rangle = n_i\,\delta_{ij}
\]
(primera ortogonalidad): $\lambda = \delta_{ij}$.

(b) Descompóngase $V$ en
\omterm{def:b3:representations:rep}{irreducibles} (Maschke): $p_i$ actúa
como la identidad sobre los sumandos de
\omterm{def:b3:representations:character}{carácter} $\chi_i$ y como $0$
sobre los demás, de modo que $p_i$ es la proyección sobre su suma $V_i$
a lo largo de la suma de los restantes; la imagen $V_i$ no depende de la
descomposición elegida (es el conjunto de vectores fijados por $p_i$,
definido sin elecciones). $\sum_ip_i$ actúa como la identidad sobre todo
sumando \omterm{def:b3:representations:rep}{irreducible}: es
$\mathrm{id}_V$. La descomposición más fina de
$V_i \cong W_i^{\oplus m_i}$ exige elegir una base de
$\operatorname{Hom}_G(W_i, V)$: canónica no es.

(c) Para $\varepsilon$ (grado $1$): $p_\varepsilon =
\frac1{6}\sum_{g}\varepsilon(g)\,\rho(g)$, es decir, en el álgebra del
grupo,
\[
p_\varepsilon = \tfrac16\bigl(e - (1\,2) - (1\,3) - (2\,3)
+ (1\,2\,3) + (1\,3\,2)\bigr) .
\]
Elevando al cuadrado: el coeficiente de $g$ en $p_\varepsilon^2$ es
$\frac1{36}\sum_{h}\varepsilon(h)\varepsilon(h^{-1}g) =
\frac1{36}\,\varepsilon(g)\sum_h\varepsilon(h)^2 =
\frac{6}{36}\varepsilon(g)$: $p_\varepsilon^2 =
p_\varepsilon$. (Su imagen en la
\omterm{ex:b3:representations:examples}{representación regular} es la
recta generada por $\sum_g\varepsilon(g)e_g$: la
\omterm{def:b3:representations:rep}{representación} signatura aparece
con multiplicidad $1$, como exige la teoría general.)
\end{solution}

\begin{solution}{exo:b3:representations:12}
(a) Ordénense las clases $e$ [1], trasposiciones [6], $3$-ciclos [8],
$4$-ciclos [6], dobles trasposiciones [3]. El segundo
\omterm{def:b3:representations:character}{carácter} lineal es
$\chi_2 = \varepsilon$, con valores $1, -1, 1, -1, 1$. Para el
\omterm{def:b3:representations:character}{carácter} $\chi_3$ de grado
$2$, la ortogonalidad de cada columna con la columna de la identidad
($\sum_in_i\chi_i(g) = 0$ para $g \neq e$) da, sobre las trasposiciones:
$1 - 1 + 2\chi_3 + 3(\chi_4 +
\chi_5) = 0$; el truco del signo $\chi_5 = \varepsilon\chi_4$ (un
\omterm{def:b3:representations:character}{carácter} de grado $3$ por uno
lineal vuelve a ser \omterm{def:b3:representations:rep}{irreducible} ---
misma norma) hace que $\chi_4 + \chi_5$ se anule en las clases impares:
$\chi_3 = 0$ allí. Sobre los $3$-ciclos: $1 + 1 +
2\chi_3(c_3) + 3(\chi_4 + \chi_5)(c_3) = 0$ con $\chi_5 = \chi_4$ en las
clases pares; la columna de $c_3$ consigo misma da $\abs{\chi_3}^2$
datos; resolviendo el pequeño sistema (úsese también la ortogonalidad de
filas de $\chi_3$ con $\mathbf 1$ y $\varepsilon$):
$\chi_3 = (2, 0, -1, 0, 2)$, y después $\chi_4 =
(3, 1, 0, -1, -1)$ y $\chi_5 = \varepsilon\chi_4 = (3, -1,
0, 1, -1)$. La tabla completa:
\begin{center}
\begin{tabular}{c|ccccc}
 & $e$ & $6\,t$ & $8\,c_3$ & $6\,c_4$ & $3\,v$\\
\hline
$\chi_1$ & $1$ & $1$ & $1$ & $1$ & $1$\\
$\chi_2$ & $1$ & $-1$ & $1$ & $-1$ & $1$\\
$\chi_3$ & $2$ & $0$ & $-1$ & $0$ & $2$\\
$\chi_4$ & $3$ & $1$ & $0$ & $-1$ & $-1$\\
$\chi_5$ & $3$ & $-1$ & $0$ & $1$ & $-1$\\
\end{tabular}
\end{center}
(Todas las filas tienen norma $1$ y todas las columnas son ortogonales:
las comprobaciones salen bien.)

(b) Los datos de clases $1, 6, 8, 6, 3$ con esos grados son los de $S_4$
(tipos de ciclo $e$, $2$, $3$, $4$, $2{+}2$). Núcleos:
$\ker\chi_2 = \{g : \varepsilon(g) = 1\} = A_4$ (clases $e, c_3, v$:
$1 + 8 + 3 = 12$, de índice $2$); $\ker\chi_3 = \{g : \chi_3(g) = 2\}$ =
clases $e, v$: el grupo de Klein $V_4$, de orden $4$, normal.
$\chi_4, \chi_5$ son fieles ($\chi_i(g) = n_i$ solo en $e$).
Intersecciones de núcleos: $\{e\}$, $V_4$, $A_4$, $G$ --- por
el~\cref{exo:b3:representations:6}(b), estos son \emph{todos} los
\omterm{def:b3:groups:normal}{subgrupos normales} de $S_4$.

(c) Los \omterm{def:b3:representations:character}{caracteres} con
$V_4 \subseteq \ker$ son $\chi_1,
\chi_2, \chi_3$: se factorizan a través de $G/V_4$, un grupo de orden
$6$ con grados \omterm{def:b3:representations:rep}{irreducibles}
$1, 1, 2$ --- la tabla de $S_3$. Como la tabla del cociente \emph{es} un
invariante completo entre los dos grupos de orden $6$ ($\Z/6\Z$ tendría
seis \omterm{def:b3:representations:character}{caracteres} lineales),
$G/V_4 \cong
S_3$: el cociente se ve dentro de la tabla como el bloque de filas que
contienen $V_4$ en su núcleo.
\end{solution}

\begin{solution}{pb:b3:representations:1}
\textbf{1.} Si $\alpha^n + c_{n-1}\alpha^{n-1} + \dots + c_0 =
0$ ($c_i \in \Z$), entonces $\alpha^n \in \Z\text{-gen}(1, \dots,
\alpha^{n-1})$ y, por inducción, toda potencia lo es: $\Z[\alpha]$ está
generado por $1, \alpha, \dots, \alpha^{n-1}$. Recíprocamente, sea
$\Z[\alpha] = \Z g_1 + \dots + \Z g_m$. Escríbase $\alpha g_i =
\sum_j m_{ij}g_j$ con $M = (m_{ij}) \in M_m(\Z)$: el vector $g = (g_i)$
cumple $(\alpha I - M)g = 0$; multiplicando por la matriz adjunta,
$\det(\alpha I - M)\,g_i = 0$ para todo $i$ y, como $1 \in \Z[\alpha]$
es una $\Z$-combinación de los $g_i$, $\det(\alpha I - M) = 0$: $\alpha$
es raíz del mónico $\det(XI - M) \in \Z[X]$.

\textbf{2.} Si $\Z[\alpha]$ está generado por las potencias de $\alpha$
hasta $n-1$ y $\Z[\beta]$ por las de $\beta$ hasta $m - 1$, entonces
$\Z[\alpha,\beta]$ está generado por los $nm$ productos
$\alpha^i\beta^j$ (redúzcase cualquier monomio). Los subanillos
$\Z[\alpha + \beta]$ y $\Z[\alpha\beta]$ son $\Z$-submódulos del
$\Z$-módulo \omterm{def:b3:modules:free}{finitamente generado}
$\Z[\alpha, \beta]$ y, por tanto,
\omterm{def:b3:modules:free}{finitamente generados} ($\Z[\alpha,\beta]$,
generado por $nm$ elementos, es imagen de $\Z^{nm}$; un submódulo se
levanta a un submódulo de $\Z^{nm}$, libre de rango $\leq nm$ por
el~\cref{thm:b3:modules:submodule}, y su imagen genera). Por la pregunta
1, $\alpha + \beta$ y $\alpha\beta$ son enteros
\omterm{def:b3:galois:algebraic}{algebraicos}.

\textbf{3.} Si $\frac pq$ (en forma \omterm{def:b3:representations:rep}{irreducible}) es raíz de un polinomio
entero mónico de grado $n$, el teorema de la raíz racional (quítense
denominadores: $p^n = -q(\cdots)$) da $q \mid p^n$, luego $q = \pm1$.

\textbf{4.} $\chi(g)$ es suma de raíces de la unidad
(la~\cref{prop:b3:representations:charbasics}), cada una de ellas entero
\omterm{def:b3:galois:algebraic}{algebraico} (raíz de $X^N - 1$);
concluir con la pregunta 2.

\textbf{5.} Para $h \in G$: $\rho(h)S_C\rho(h)^{-1} = \sum_{g\in
C}\rho(hgh^{-1}) = S_C$ ($C$ es una clase). Por Schur, $S_C =
\omega_C\,\mathrm{id}$; tomando trazas, $\abs C\,\chi(g_C) =
\omega_C\, n$.

\textbf{6.} $S_CS_{C'} = \sum_{x \in C, y \in C'}\rho(xy) =
\sum_{z \in G} a(z)\rho(z)$ con $a(z) = \#\{(x,y) \in C\times
C' : xy = z\}$. La conjugación por $h$ biyecta las soluciones para $z$
con las de $hzh^{-1}$: $a$ es una
\omterm{def:b3:representations:character}{función de clase} con valores
en $\N$, luego $S_CS_{C'} = \sum_{C''}a_{CC'C''}S_{C''}$. Sustituyendo
$S_C = \omega_C\,\mathrm{id}$ en todo e identificando los escalares:
$\omega_C\omega_{C'} = \sum_{C''}a_{CC'C''}\,\omega_{C''}$.

\textbf{7.} Sea $M$ el $\Z$-módulo generado por $1$ y por todos los
productos $\omega_{C_1}\cdots\omega_{C_k}$; por la pregunta 6, todo
producto de ese tipo se reduce a una $\Z$-combinación de $1$ y de los
$\omega_C$: $M$ es \omterm{def:b3:modules:free}{finitamente generado}, y
$\omega_C M
\subseteq M$ para cada $C$. En particular, $\Z[\omega_C]
\subseteq M$ es \omterm{def:b3:modules:free}{finitamente generado}
(submódulo, como en la pregunta 2), y la pregunta 1 hace de $\omega_C$
un entero \omterm{def:b3:galois:algebraic}{algebraico}.

\textbf{8.} Para una $\chi$
\omterm{def:b3:representations:rep}{irreducible} de grado $n$:
\[
\frac{\abs G}{n} = \frac{\abs G}{n}\langle\chi,\chi\rangle
= \frac1n \sum_{g}\chi(g)\overline{\chi(g)}
= \sum_{C}\frac{\abs C\,\chi(g_C)}{n}\,\overline{\chi(g_C)}
= \sum_C \omega_C\,\overline{\chi(g_C)} .
\]
Cada $\overline{\chi(g_C)} = \chi(g_C^{-1})$ es un entero
\omterm{def:b3:galois:algebraic}{algebraico} (pregunta 4), luego el
miembro de la derecha lo es (preguntas 2 y 7); y es racional, luego
entero (pregunta 3): $n \mid \abs
G$.

\textbf{9.} Bézout: $u\abs C + vn = 1$ con $u, v \in \Z$. Entonces
\[
\frac{\chi(g_C)}{n} = u\,\frac{\abs C\,\chi(g_C)}{n} +
v\,\chi(g_C) = u\,\omega_C + v\,\chi(g_C),
\]
un entero \omterm{def:b3:galois:algebraic}{algebraico}.

\textbf{10.} Supongamos $0 < \abs{\chi(g_C)} < n$ y sea $\alpha
= \chi(g_C)/n$. Todos los valores $\chi(g)$ están en $\Q(\zeta_N)$, $N =
\abs G$ (sumas de raíces $N$-ésimas de la unidad). Para $\sigma \in
\operatorname{Gal}(\Q(\zeta_N)/\Q)$: $\sigma$ lleva raíces de la unidad
a raíces de la unidad ($\sigma(\zeta^k) = \zeta^{ak}$,
el~\cref{thm:b3:galois:cyclotomicirred}), de modo que
$\sigma(\chi(g_C))$ vuelve a ser suma de $n$ raíces de la unidad:
$\abs{\sigma(\alpha)}
\leq 1$; además, $\sigma(\alpha)$ es un entero
\omterm{def:b3:galois:algebraic}{algebraico} (tiene el mismo
\omterm{def:b3:galois:algebraic}{polinomio mínimo} que $\alpha$). El
producto $P =
\prod_\sigma \sigma(\alpha)$ queda fijado por todo el
\omterm{def:b3:galois:galois}{grupo de Galois}, luego es racional
(el~\cref{thm:b3:galois:fundamental}(1)), y es un entero
\omterm{def:b3:galois:algebraic}{algebraico} con
\[
0 < \abs P \leq \abs\alpha < 1
\]
(ningún factor se anula: $\sigma(\alpha) = 0$ obligaría a $\alpha =
0$). Esto contradice la pregunta 3. Por tanto, $\chi(g_C) = 0$ o
$\abs{\chi(g_C)} = n$ y, en el segundo caso, $\rho(g_C)$ es escalar
(la~\cref{prop:b3:representations:charbasics}).

\textbf{11.} Ortogonalidad de columnas ($C \neq \{e\}$):
$\sum_i \chi_i(e)\overline{\chi_i(g_C)} = 0$, es decir, $1 +
\sum_{i \geq 2} n_i\overline{\chi_i(g_C)} = 0$. Si todo $\chi_i$ no
trivial con $p \nmid n_i$ se anulara en $g_C$, agrupando el resto por su
factor $p$:
\[
-\frac1p = \sum_{i \geq 2,\ p \mid n_i}
\frac{n_i}{p}\,\overline{\chi_i(g_C)},
\]
un entero \omterm{def:b3:galois:algebraic}{algebraico} --- en contra de
la pregunta 3. Luego algún $\chi_i$ no trivial cumple $p \nmid n_i$ y
$\chi_i(g_C) \neq 0$; como $\abs C = p^k$, $\gcd(\abs C, n_i) = 1$, y la
pregunta 10 hace de $\rho_i(g_C)$ un escalar. Ahora bien, $G$ es simple
no abeliano: $\chi_i$ es fiel (el~\cref{exo:b3:representations:6}(c)) y
$Z_i = \{g : \rho_i(g) \text{ escalar}\}$ es un
\omterm{def:b3:groups:normal}{subgrupo normal} (la preimagen por
$\rho_i$ de los escalares, que forman un \omterm{def:b3:groups:normal}{subgrupo normal} --- de hecho
central --- de la imagen) que contiene $g_C \neq e$: $Z_i = G$. Entonces
$\rho_i(G)$ es abeliano y fiel, lo que hace $G$ abeliano: contradicción.
\emph{Ningún \omterm{def:b3:groups:simple}{grupo simple} no abeliano tiene una
\omterm{ex:b3:groups:actions}{clase de conjugación} de cardinal potencia
de primo
$> 1$.}

\textbf{12.} Sea $G$ simple de orden $p^aq^b$. Si $b = 0$ ($G$ es un
$p$-grupo): $Z(G) \neq \{e\}$ (el~\cref{thm:b3:groups:pfixed}) es
normal, luego $Z(G) = G$: abeliano. En caso contrario, tómense $Q$ un
$q$-subgrupo de Sylow y $z \in Z(Q)
\setminus\{e\}$ (de nuevo el~\cref{thm:b3:groups:pfixed}). Entonces $Q
\subseteq Z_G(z)$, de modo que la clase de $z$ tiene cardinal
$[G : Z_G(z)]$ divisor de $[G : Q] = p^a$: una potencia de $p$. Si el
cardinal es $1$, $z \in Z(G)$: el \omterm{ex:b3:groups:actions}{centro}
es un \omterm{def:b3:groups:normal}{subgrupo normal} no trivial, luego
$Z(G) = G$, abeliano. Si el cardinal es $p^k > 1$: la pregunta 11 lo
prohíbe para $G$ simple no abeliano. En cualquier caso, un
\omterm{def:b3:groups:simple}{grupo simple} de orden $p^aq^b$ es
abeliano ($\cong \Z/p\Z$).

\textbf{13.} Inducción sobre $\abs G$ ($\abs G = 1$:
\omterm{def:b3:groups:derived}{resoluble}). Si $G$ es simple, la
pregunta 12 lo hace abeliano y, por tanto,
\omterm{def:b3:groups:derived}{resoluble}. En caso contrario, tómese $N$
un \omterm{def:b3:groups:normal}{subgrupo normal} propio no trivial:
$\abs
N$ y $\abs{G/N}$ son de nuevo de la forma $p^{a'}q^{b'}$ y más pequeños,
de modo que $N$ y $G/N$ son \omterm{def:b3:groups:derived}{resolubles}
por inducción, y $G$ es \omterm{def:b3:groups:derived}{resoluble}
(la~\cref{prop:b3:groups:derived}). --- Con tres primos falla el paso
clave: el índice de un \omterm{def:b3:groups:sylow}{subgrupo de Sylow}
ya no es una potencia de primo, de modo que la clase de un elemento
central de un Sylow no tiene por qué tener cardinal potencia de primo. Y
algún fallo es inevitable: $A_5$, de orden $2^2\cdot3\cdot5$, es simple
y no \omterm{def:b3:groups:derived}{resoluble}.

\textbf{14.} Las clases y sus cardinales son los
del~\cref{exo:b3:groups:11}(a). La $S_5$-clase de $c$ tiene cardinal
$24$; si siguiera siendo una sola $A_5$-clase, el recuento
órbita--estabilizador daría $\abs{Z_{A_5}(c)} = 60/24$, que no es entero
--- concretamente, $Z_{S_5}(c) = \langle c\rangle$ tiene orden $5$ y
está dentro de $A_5$, de modo que la $A_5$-clase de $c$ tiene cardinal
$60/5 =
12$: la $S_5$-clase se desdobla en dos ($c$ y $c^2$ solo son conjugados
en $S_5$ mediante una permutación impar).

\textbf{15.} Un solo \omterm{def:b3:representations:character}{carácter}
lineal: una \omterm{def:b3:representations:rep}{representación} de grado
$1$ se factoriza a través de $G/D(G)$, y $D(A_5) = A_5$ ($A_5$ es simple
no abeliano y $D(A_5)$ es normal y no trivial --- $A_5$ no es abeliano).
Luego $n_1 = 1$ y $n_2^2 + n_3^2 + n_4^2
+ n_5^2 = 59$ con cada $n_i \geq 2$. Cuadrados disponibles: $4, 9,
16, 25, 36, 49$. Una suma de cuatro de ellos igual a $59$: el mayor ha
de ser $25$ ($36 + 4 + 4 + 9 = 53 < 59$ no permite ajustar:
$36 + 16 + 4 + 4 = 60$, $36 + 9 + 9 + 4 = 58$, $36 +
16 + 9 + 4 = 65$ --- ninguna combinación con $36$ o $49$ funciona), y
$59 - 25 = 34 = 16 + 9 + 9$ (la única manera: $16 + 16 + 4 =
36$, $25 + 9 + 4 = 38$, $25 + 4 + 4 = 33$): grados $1, 3, 3,
4, 5$.

\textbf{16.} $\langle\pi, \mathbf 1\rangle = \frac1{60}(5 +
15\cdot1 + 20\cdot2 + 0 + 0) = 1$ (una sola
\omterm{def:b3:groups:action}{órbita} --- recuento de Burnside) y
$\langle\pi, \pi\rangle = \frac1{60}(25 + 15 + 80 +
0 + 0) = 2$: $\pi$ contiene el
\omterm{def:b3:representations:character}{carácter} trivial una vez, y
su otro constituyente es un único
\omterm{def:b3:representations:rep}{irreducible}. Por tanto,
$\chi_4 = \pi
- \mathbf 1$ es \omterm{def:b3:representations:rep}{irreducible} de
grado $4$, con valores $4, 0,
1, -1, -1$.

\textbf{17.} Sobre los pares, un elemento fija $\{i,j\}$ si y solo si
fija o intercambia $i, j$: los recuentos son $10$ ($e$), $2$ ($t =
(1\,2)(3\,4)$ fija $\{1,2\}, \{3,4\}$), $1$ ($(1\,2\,3)$ fija
$\{4,5\}$), $0, 0$. Entonces $\langle\pi_{10},
\pi_{10}\rangle = \frac1{60}(100 + 15\cdot4 + 20\cdot1) = 3$: tres
constituyentes \omterm{def:b3:representations:rep}{irreducibles}, cada
uno una vez. Y
$\langle\pi_{10}, \mathbf 1\rangle = \frac1{60}(10 + 30 + 20)
= 1$, $\langle\pi_{10}, \chi_4\rangle = \frac1{60}(40 + 0 +
20\cdot1\cdot1 + 0 + 0) = 1$: luego $\pi_{10} = \mathbf 1 + \chi_4
+ \chi$ con $\chi$ \omterm{def:b3:representations:rep}{irreducible} de
grado $10 - 1 - 4 = 5$ y valores
$\chi_5 = \pi_{10} - \mathbf 1 - \chi_4 = (5, 1, -1, 0,
0)$.

\textbf{18.} Escribamos $a, a'$ para los valores de $\chi_2, \chi_3$ en
$t$, y $b, b'$ para los de $s = (1\,2\,3)$; los cuatro son reales ($t$ y
$s$ son conjugados de sus inversos). Columna $t$ contra columna $e$:
$1 + 3a + 3a' + 4\cdot0 + 5\cdot1 = 0$, luego $a + a'
= -2$; columna $t$ consigo misma: $1 + a^2 + a'^2 + 0 + 1 =
\frac{60}{15} = 4$, luego $a^2 + a'^2 = 2$; de donde $(a + a')^2 = 4
= 2 + 2aa'$ da $aa' = 1$ y $a = a' = -1$. Columna $s$ contra $e$:
$1 + 3(b + b') + 4 - 5 = 0$ da $b + b' = 0$; columna $s$ consigo misma:
$1 + b^2 + b'^2 + 1 + 1 = \frac{60}
{20} = 3$ da $b = b' = 0$.

\textbf{19.} Columna $c$ contra columna $e$: $1 + 3(x + y) +
4(-1) + 5\cdot0 = 0$, luego $x + y = 1$. Columna $c$ contra columna
$c^2$ (clases distintas, ortogonales): $1 + xy + yx + 1
+ 0 = 0$, luego $xy = -1$. Así, $x, y$ resuelven $T^2 - T - 1 = 0$:
$\{x, y\} = \{\varphi, \bar\varphi\}$ con $\varphi = \frac{1 +
\sqrt5}2$. Coherencia: $x^2 + y^2 = (x+y)^2 - 2xy = 3 =
\frac{60}{12} - 2$ --- lo que concuerda con la identidad de la columna
consigo misma $1 +
x^2 + y^2 + 1 + 0 = 5$. La tabla:
\begin{center}
\begin{tabular}{c|ccccc}
 & $e$ & $15\,t$ & $20\,s$ & $12\,c$ & $12\,c^2$\\
\hline
$\chi_1$ & $1$ & $1$ & $1$ & $1$ & $1$\\
$\chi_2$ & $3$ & $-1$ & $0$ & $\varphi$ & $\bar\varphi$\\
$\chi_3$ & $3$ & $-1$ & $0$ & $\bar\varphi$ & $\varphi$\\
$\chi_4$ & $4$ & $0$ & $1$ & $-1$ & $-1$\\
$\chi_5$ & $5$ & $1$ & $-1$ & $0$ & $0$\\
\end{tabular}
\end{center}

\textbf{20.} $\norm{\chi_2}^2 = \frac1{60}\bigl(9 + 15\cdot1 +
0 + 12\varphi^2 + 12\bar\varphi^2\bigr) = \frac{9 + 15 +
36}{60} = 1$, usando $\varphi^2 + \bar\varphi^2 = (\varphi +
\bar\varphi)^2 - 2\varphi\bar\varphi = 1 + 2 = 3$. Análogamente,
$\langle\chi_2, \chi_3\rangle = \frac1{60}(9 + 15 + 0 +
12(2\varphi\bar\varphi)) = \frac{9 + 15 - 24}{60} = 0$.
Divisibilidad de \omterm{thm:b3:galois:finitefields}{Frobenius}:
$1, 3, 3, 4, 5$ dividen todos a $60$. Los valores irracionales
$\varphi, \bar\varphi$ son la diferencia visible con $S_5$: en $S_5$
todo elemento es conjugado de todos los generadores de su grupo cíclico
con el mismo tipo de ciclo --- en particular $c \sim c^2$ ---, lo que
fuerza valores \omterm{def:b3:representations:character}{carácter}
racionales (de hecho enteros); en $A_5$, el desdoblamiento de los ciclos
de longitud cinco abre la puerta a $\Q(\sqrt5)$.

\textbf{21.} Un \omterm{def:b3:groups:normal}{subgrupo normal} $N$ es
unión de clases, contiene $e$ y $\abs N \mid 60$. Las sumas candidatas:
$1{+}15 = 16$, $1{+}20 = 21$, $1{+}12 = 13$, $1{+}24 = 25$,
$1{+}15{+}20 = 36$, $1{+}15{+}12 = 28$, $1{+}15{+}24 = 40$,
$1{+}20{+}12 = 33$, $1{+}20{+}24 = 45$, $1{+}12{+}12 = 25$,
$1{+}15{+}20{+}12 = 48$, $1{+}15{+}20{+}24 = 60 - 12 = \dots$
enumerando todas las subsumas propias que contienen $1$: ninguna de
$13, 16,
21, 25, 28, 33, 36, 40, 45, 48, 13{+}\dots$ divide a $60$ salvo el
propio $1$: $N = \{e\}$ o $A_5$. Simplicidad, leída en cinco números. Y
el criterio de la pregunta 11 también se ve: los cardinales de clase
$15 = 3\cdot5$, $20 = 4\cdot5$, $12 = 4\cdot3$ son todos compuestos de
dos primos --- ninguna clase de cardinal potencia de primo, exactamente
como el criterio de Burnside exige de un
\omterm{def:b3:groups:simple}{grupo simple}.

\textbf{22.} Una rotación de $\R^3$ de ángulo $\theta$ tiene valores
propios $1, \eu^{\iu\theta}, \eu^{-\iu\theta}$: traza $1 +
2\cos\theta$. Para $\theta = \frac{2\pi}5$: $2\cos\frac{2\pi}5
= \frac{\sqrt5 - 1}2$, luego la traza es $1 + \frac{\sqrt5-1}2 =
\frac{1+\sqrt5}2 = \varphi$. El elemento no trivial $\tau$ de
$\operatorname{Gal}(\Q(\sqrt5)/\Q)$, aplicado entrada a entrada a una
\omterm{def:b3:representations:table}{tabla de caracteres}, lleva
\omterm{def:b3:representations:character}{caracteres} a
\omterm{def:b3:representations:character}{caracteres} (conmuta con el
álgebra que los define: $\tau\circ\chi$ es el
\omterm{def:b3:representations:character}{carácter} de la
\omterm{def:b3:representations:rep}{representación} obtenida
transportando las matrices mediante $\tau$ sobre las entradas; o, en
abstracto: las relaciones de ortogonalidad son $\Q$-racionales, de modo
que $\tau$ permuta sus soluciones); $\tau$ fija $\chi_1, \chi_4, \chi_5$
(valores racionales) y, por tanto, ha de intercambiar $\chi_2$ y
$\chi_3$: el segundo \omterm{def:b3:representations:character}{carácter}
de grado $3$ lleva los valores conjugados, sin calcular ninguna matriz.
Geométricamente, las dos
representaciones son la acción
icosaédrica y su composición con un automorfismo exterior de $A_5$ (la
conjugación por una trasposición), que intercambia las dos clases de
ciclos de longitud cinco.

\textbf{23.} Para $z \in Z(G)$, $\rho(z)$ conmuta con todo $\rho(g)$, de
modo que, por el lema de Schur, $\rho(z) = \lambda\,
\mathrm{id}$; y como $z$ tiene orden finito, $\lambda$ es una raíz de la
unidad y $\abs{\chi(z)} = \abs\lambda\,n = n$. Entonces
\[
\abs G = \abs G\,\langle\chi, \chi\rangle
= \sum_{g \in G}\abs{\chi(g)}^2
\geq \sum_{z \in Z(G)}\abs{\chi(z)}^2
= \abs{Z(G)}\,n^2,
\]
es decir, $n^2 \leq [G : Z(G)]$. Un grupo no abeliano de orden $8$
($D_4$ o $Q_8$) tiene cinco \omterm{ex:b3:groups:actions}{clases de
conjugación}, luego cinco grados
\omterm{def:b3:representations:rep}{irreducibles} con $\sum n_i^2 = 8$;
la única manera de escribir $8$ como suma de cinco cuadrados $\geq 1$ es
$1 + 1 + 1 + 1
+ 4$: grados $1, 1, 1, 1, 2$. Ambos grupos tienen
\omterm{ex:b3:groups:actions}{centro} de orden $2$, y el
\omterm{def:b3:representations:character}{carácter} de grado $2$ alcanza
la igualdad: $2^2 = 4 = [G : Z(G)]$ --- la cota es óptima. Para $A_5$,
$Z = \{e\}$ y la cota se lee $n^2 \leq 60$: se cumple con
$1, 3, 3, 4, 5$ y con holgura ($25 \leq 60$), como no podía ser de otro
modo, pues la igualdad forzaría (por la misma cadena) que $\chi$ se
anulase fuera del \omterm{ex:b3:groups:actions}{centro}.

\textbf{24.} $\rho(g)^m = \mathrm{id}$, siendo $m$ el orden de $g$, de
modo que $\rho(g)$ es aniquilado por $X^m - 1$, que se descompone con
raíces simples sobre $\C$: es diagonalizable, con valores propios
$\lambda_1, \dots, \lambda_n$ raíces de la unidad, en una base propia
$(e_i)$. Los productos $e_ie_j$ ($i \leq j$) forman una base propia de
$\operatorname{Sym}^2V$ con valores propios $\lambda_i\lambda_j$, y los
$e_i \wedge e_j$ ($i < j$) una de $\Lambda^2V$; puesto que
\[
\sum_{i<j}\lambda_i\lambda_j
= \frac{(\sum_i\lambda_i)^2 - \sum_i\lambda_i^2}2
= \frac{\chi(g)^2 - \chi(g^2)}2,
\]
y la suma simétrica añade $\sum_i\lambda_i^2$ en lugar de restarlo,
ambas fórmulas se siguen. Para $\chi_2 = (3, -1,
0, \varphi, \bar\varphi)$ sobre las clases $(e, (2,2)\text{-},
3\text{-ciclos}, c, c^2)$: elevar al cuadrado envía las dobles
trasposiciones a $e$, los ciclos de longitud tres a ciclos de longitud
tres, y la clase de $c$ sobre la de $c^2$ y recíprocamente ($c^{-1} \sim
c$ en $A_5$, luego $c^4 \sim c$). Por tanto, $\chi_2(g^2)$ se lee $(3,
3, 0, \bar\varphi, \varphi)$ y, usando $\varphi^2 = \varphi
+ 1$, $\bar\varphi = 1 - \varphi$:
\[
\Lambda^2\chi_2 = (3, -1, 0, \varphi, \bar\varphi) = \chi_2,
\qquad
\operatorname{Sym}^2\chi_2 = (6, 2, 0, 1, 1) .
\]
Descomponiendo esta última, con los cardinales de clase $1, 15, 20, 12,
12$: $\langle\cdot, \mathbf 1\rangle = \frac1{60}(6 +
15\cdot2 + 0 + 12 + 12) = 1$; $\langle\cdot, \chi_5\rangle =
\frac1{60}(30 + 30) = 1$; $\langle\cdot, \chi_4\rangle =
\frac1{60}(24 - 12 - 12) = 0$; $\langle\cdot, \chi_2\rangle =
\frac1{60}(18 - 30 + 12(\varphi + \bar\varphi)) = 0$, y análogamente
para $\chi_3$. Luego $\operatorname{Sym}^2\chi_2 =
\mathbf 1 + \chi_5$ (dimensiones $6 = 1 + 5$) y
$\chi_2\otimes\chi_2 = \mathbf 1 + \chi_2 + \chi_5$
(dimensiones $9 = 1 + 3 + 5$). Geometría: el isomorfismo equivariante
$\Lambda^2\R^3 \to \R^3$, $u \wedge v \mapsto u
\times v$, es exactamente $\Lambda^2\chi_2 = \chi_2$ para un grupo de
rotaciones; el sumando $\mathbf 1$ del cuadrado simétrico es la forma
cuadrática invariante $x^2 + y^2 + z^2$, y $\chi_5$ vive en el espacio
de dimensión cinco de los tensores simétricos de traza nula (las
cuádricas armónicas).

\textbf{25.} Columna de $c$ contra columna de $c^2$:
\[
1\cdot1 + \varphi\bar\varphi + \bar\varphi\varphi +
(-1)(-1) + 0 = 1 - 1 - 1 + 1 + 0 = 0,
\]
tal como exige la ortogonalidad para clases distintas ($\varphi
\bar\varphi = -1$). Columna de $c$ consigo misma: $1 +
\varphi^2 + \bar\varphi^2 + 1 + 0 = 1 + 3 + 1 = 5 = 60/12 =
\abs{Z_{A_5}(c)}$. \omterm{def:b3:representations:character}{Carácter}
regular $\sum_in_i\chi_i$ en las cuatro columnas distintas de la
identidad:
\[
1 - 3 - 3 + 0 + 5 = 0, \qquad
1 + 0 + 0 + 4 - 5 = 0,
\]
sobre las dobles trasposiciones y los ciclos de longitud tres, y sobre
la clase de $c$ (la de $c^2$ es su conjugada
\omterm{def:b3:galois:galois}{de Galois}):
\[
1 + 3\varphi + 3\bar\varphi - 4 + 0 = 1 + 3 - 4 = 0,
\]
usando $\varphi + \bar\varphi = 1$. La tabla supera todas las
auditorías: es la \omterm{def:b3:representations:table}{tabla de
caracteres} de $A_5$.
\end{solution}
